% --- Archon physics formalization source begin ---
% archon:physics
% archon:covers PhyXMiniProblems/problem_phyx_mini_0632.lean
% archon:source-report reports/phyx_mini/problem_phyx_mini_0632.source.json

\chapter{Physics problem phyx\_mini\_0632}
\label{ch:PhyXMiniProblems_problem_phyx_mini_0632}

\paragraph{Problem source.}
\#\# Physical scenario

Consider the following hypothetical universe: The Nibiruvians and the Xibalbans are minuscule “flat” societies located at an angular separation of \$\textbackslash{}theta = 60\textasciicircum{}\{\textbackslash{}circ\}\$ on the surface of a two-dimensional spherical universe with radius \$R\$, as shown in figure. The radius of this universe was \$R\_0 = 500.0\textbackslash{} m\$ upon its creation at time \$t=0\$, and it has been increasing at a constant rate of \$1.00\textbackslash{} \textbackslash{}mu m/s\$.

\#\# Question

What is the distance between Nibiru and Xibalba four years after creation?

\#\# Answer choices

- A: A: 656m

- B: B: 626m

- C: C: 636m

- D: D: 646m

\#\# Figure caption (auxiliary; use the image as primary evidence)

The image is a diagram featuring a circle with a blue gradient. It includes two labeled points on the circumference named "Xibalba" and "Nibiru." A curved magenta line connects these two points, labeled \textbackslash{}(D(t)\textbackslash{}).

There are two dashed black lines:

1. A horizontal line from "Xibalba" to the center
2. A vertical line from the center to "Nibiru"

These lines form a right triangle with the center of the circle. The angle between the horizontal line and the hypotenuse is labeled \textbackslash{}(\textbackslash{}theta\textbackslash{}).

The hypotenuse of the triangle, extending from the center to "Nibiru," is labeled \textbackslash{}(R(t)\textbackslash{}). Both the horizontal and vertical dashed lines have arrowheads indicating direction toward the center and "Nibiru," respectively.

\#\# Dataset metadata

Relativity; Physical Model Grounding Reasoning, Spatial Relation Reasoning

\paragraph{Recorded answer/context.}
A: A: 656m

\paragraph{Figure/image path.}
/root/proposal\_for\_physic/science-mango/phyx\_mini\_run/phyx\_data/test\_image/632.png

\paragraph{Formalization target.}
create a compiling Lean file with sorry bodies at `PhyXMiniProblems/problem\_phyx\_mini\_0632.lean`. The Lean declarations must preserve the physical quantities, dimensions or dimensional roles, figure labels, governing-law hypotheses, and final relation expressed by this problem.
Use Mathlib/Physlib names found through LeanExplore where available. If a physics API is missing, introduce faithful local abstractions rather than scalar placeholder aliases.

\begin{theorem}[Physics formalization target]
\label{thm:physics:phyx_mini_0632:target}
\lean{PhyXMiniProblems.ProblemPhyXMini0632.problem_phyx_mini_0632}
\uses{def:physics:phyx-mini-0632:phyxminiproblems-problemphyxmini0632-expandingsphericaluniversesetup, def:physics:phyx-mini-0632:phyxminiproblems-problemphyxmini0632-matchesexpandinguniverseproblemdata, def:physics:phyx-mini-0632:phyxminiproblems-problemphyxmini0632-matchessuppliedexpandinguniversefigure, def:physics:phyx-mini-0632:phyxminiproblems-problemphyxmini0632-hasphysicalexpandinguniverseparameters, def:physics:phyx-mini-0632:phyxminiproblems-problemphyxmini0632-satisfiesconstantexpansionandsphericalarclaws, def:physics:phyx-mini-0632:phyxminiproblems-problemphyxmini0632-recordeddatasetanswer, def:physics:phyx-mini-0632:phyxminiproblems-problemphyxmini0632-roundstonearestmeter, def:physics:phyx-mini-0632:phyxminiproblems-problemphyxmini0632-isuniquematchinganswerchoice}
The assigned autoformalize agent should translate this physics problem into Lean declarations in the covered file, with theorem and lemma proof bodies written as `by sorry`.
\end{theorem}
\begin{proof}
This is an autoformalization task, not a proof task. Produce statements that can later be proved by the physics prover without weakening the physical model.
\end{proof}

% --- Archon declaration topology begin ---
\section{Declaration topology}
\begin{definition}[Length Quantity]
  \label{def:physics:phyx-mini-0632:phyxminiproblems-problemphyxmini0632-lengthquantity}
  \lean{PhyXMiniProblems.ProblemPhyXMini0632.LengthQuantity}
  A nonnegative, unit-independent physical length
\end{definition}
\begin{proof}
  This is the declared model definition of Length Quantity.
\end{proof}
\begin{definition}[Time Quantity]
  \label{def:physics:phyx-mini-0632:phyxminiproblems-problemphyxmini0632-timequantity}
  \lean{PhyXMiniProblems.ProblemPhyXMini0632.TimeQuantity}
  A nonnegative, unit-independent physical time
\end{definition}
\begin{proof}
  This is the declared model definition of Time Quantity.
\end{proof}
\begin{definition}[length Readout]
  \label{def:physics:phyx-mini-0632:phyxminiproblems-problemphyxmini0632-lengthreadout}
  \lean{PhyXMiniProblems.ProblemPhyXMini0632.lengthReadout}
  \uses{def:physics:phyx-mini-0632:phyxminiproblems-problemphyxmini0632-lengthquantity}
  Read a physical length in a selected Physlib length unit
\end{definition}
\begin{proof}
  This is the declared model definition of length Readout.
\end{proof}
\begin{definition}[time Readout]
  \label{def:physics:phyx-mini-0632:phyxminiproblems-problemphyxmini0632-timereadout}
  \lean{PhyXMiniProblems.ProblemPhyXMini0632.timeReadout}
  \uses{def:physics:phyx-mini-0632:phyxminiproblems-problemphyxmini0632-timequantity}
  Read a physical time in a selected Physlib time unit
\end{definition}
\begin{proof}
  This is the declared model definition of time Readout.
\end{proof}
\begin{definition}[speed Readout]
  \label{def:physics:phyx-mini-0632:phyxminiproblems-problemphyxmini0632-speedreadout}
  \lean{PhyXMiniProblems.ProblemPhyXMini0632.speedReadout}
  Read a physical speed in selected compatible length and time units
\end{definition}
\begin{proof}
  This is the declared model definition of speed Readout.
\end{proof}
\begin{definition}[length In Meters]
  \label{def:physics:phyx-mini-0632:phyxminiproblems-problemphyxmini0632-lengthinmeters}
  \lean{PhyXMiniProblems.ProblemPhyXMini0632.lengthInMeters}
  \uses{def:physics:phyx-mini-0632:phyxminiproblems-problemphyxmini0632-lengthquantity, def:physics:phyx-mini-0632:phyxminiproblems-problemphyxmini0632-lengthreadout}
  Metre readout of a physical length
\end{definition}
\begin{proof}
  This is the declared model definition of length In Meters.
\end{proof}
\begin{definition}[time In Seconds]
  \label{def:physics:phyx-mini-0632:phyxminiproblems-problemphyxmini0632-timeinseconds}
  \lean{PhyXMiniProblems.ProblemPhyXMini0632.timeInSeconds}
  \uses{def:physics:phyx-mini-0632:phyxminiproblems-problemphyxmini0632-timequantity, def:physics:phyx-mini-0632:phyxminiproblems-problemphyxmini0632-timereadout}
  Second readout of a physical time
\end{definition}
\begin{proof}
  This is the declared model definition of time In Seconds.
\end{proof}
\begin{definition}[speed In Meters Per Second]
  \label{def:physics:phyx-mini-0632:phyxminiproblems-problemphyxmini0632-speedinmeterspersecond}
  \lean{PhyXMiniProblems.ProblemPhyXMini0632.speedInMetersPerSecond}
  \uses{def:physics:phyx-mini-0632:phyxminiproblems-problemphyxmini0632-speedreadout}
  Metres-per-second readout of a physical speed
\end{definition}
\begin{proof}
  This is the declared model definition of speed In Meters Per Second.
\end{proof}
\begin{definition}[problem Year]
  \label{def:physics:phyx-mini-0632:phyxminiproblems-problemphyxmini0632-problemyear}
  \lean{PhyXMiniProblems.ProblemPhyXMini0632.problemYear}
  The multiple-choice calculation uses the elementary convention that one problem year has 365 twenty-four-hour days. This is a unit convention, not an assumption about the requested distance
\end{definition}
\begin{proof}
  This is the declared model definition of problem Year.
\end{proof}
\begin{definition}[four Problem Years In Seconds]
  \label{def:physics:phyx-mini-0632:phyxminiproblems-problemphyxmini0632-fourproblemyearsinseconds}
  \lean{PhyXMiniProblems.ProblemPhyXMini0632.fourProblemYearsInSeconds}
  The scalar number of seconds in the four problem years
\end{definition}
\begin{proof}
  This is the declared model definition of four Problem Years In Seconds.
\end{proof}
\begin{definition}[Society]
  \label{def:physics:phyx-mini-0632:phyxminiproblems-problemphyxmini0632-society}
  \lean{PhyXMiniProblems.ProblemPhyXMini0632.Society}
  The two societies named on the boundary of the spherical universe
\end{definition}
\begin{proof}
  This is the declared model definition of Society.
\end{proof}
\begin{definition}[Figure Quantity Label]
  \label{def:physics:phyx-mini-0632:phyxminiproblems-problemphyxmini0632-figurequantitylabel}
  \lean{PhyXMiniProblems.ProblemPhyXMini0632.FigureQuantityLabel}
  Literal symbolic labels visible in the supplied figure
\end{definition}
\begin{proof}
  This is the declared model definition of Figure Quantity Label.
\end{proof}
\begin{definition}[Figure Quantity Role]
  \label{def:physics:phyx-mini-0632:phyxminiproblems-problemphyxmini0632-figurequantityrole}
  \lean{PhyXMiniProblems.ProblemPhyXMini0632.FigureQuantityRole}
  Physical role denoted by each symbolic label in the figure
\end{definition}
\begin{proof}
  This is the declared model definition of Figure Quantity Role.
\end{proof}
\begin{definition}[Expanding Spherical Universe Figure]
  \label{def:physics:phyx-mini-0632:phyxminiproblems-problemphyxmini0632-expandingsphericaluniversefigure}
  \lean{PhyXMiniProblems.ProblemPhyXMini0632.ExpandingSphericalUniverseFigure}
  \uses{def:physics:phyx-mini-0632:phyxminiproblems-problemphyxmini0632-society, def:physics:phyx-mini-0632:phyxminiproblems-problemphyxmini0632-figurequantitylabel, def:physics:phyx-mini-0632:phyxminiproblems-problemphyxmini0632-figurequantityrole}
  Qualitative evidence transcribed from image 632.png. The raster is schematic: it gives topology, labels, and the central-angle construction but no numerical spatial scale
\end{definition}
\begin{proof}
  This is the declared model definition of Expanding Spherical Universe Figure.
\end{proof}
\begin{definition}[Expanding Spherical Universe Setup]
  \label{def:physics:phyx-mini-0632:phyxminiproblems-problemphyxmini0632-expandingsphericaluniversesetup}
  \lean{PhyXMiniProblems.ProblemPhyXMini0632.ExpandingSphericalUniverseSetup}
  \uses{def:physics:phyx-mini-0632:phyxminiproblems-problemphyxmini0632-lengthquantity, def:physics:phyx-mini-0632:phyxminiproblems-problemphyxmini0632-timequantity, def:physics:phyx-mini-0632:phyxminiproblems-problemphyxmini0632-expandingsphericaluniversefigure}
  The time-dependent radius R(t) and the Nibiru--Xibalba surface distance D(t) are independent observables. The governing laws below relate them
\end{definition}
\begin{proof}
  This is the declared model definition of Expanding Spherical Universe Setup.
\end{proof}
\begin{definition}[Matches Expanding Universe Problem Data]
  \label{def:physics:phyx-mini-0632:phyxminiproblems-problemphyxmini0632-matchesexpandinguniverseproblemdata}
  \lean{PhyXMiniProblems.ProblemPhyXMini0632.MatchesExpandingUniverseProblemData}
  \uses{def:physics:phyx-mini-0632:phyxminiproblems-problemphyxmini0632-timereadout, def:physics:phyx-mini-0632:phyxminiproblems-problemphyxmini0632-speedreadout, def:physics:phyx-mini-0632:phyxminiproblems-problemphyxmini0632-lengthinmeters, def:physics:phyx-mini-0632:phyxminiproblems-problemphyxmini0632-timeinseconds, def:physics:phyx-mini-0632:phyxminiproblems-problemphyxmini0632-problemyear, def:physics:phyx-mini-0632:phyxminiproblems-problemphyxmini0632-expandingsphericaluniversesetup}
  Numerical readouts stated in the problem: a 500 m initial radius, constant rate magnitude 1 micrometre/second, a four-year observation time, and a 60 degree = pi/3 radian central separation. No distance answer occurs here
\end{definition}
\begin{proof}
  This is the declared model definition of Matches Expanding Universe Problem Data.
\end{proof}
\begin{definition}[Matches Supplied Expanding Universe Figure]
  \label{def:physics:phyx-mini-0632:phyxminiproblems-problemphyxmini0632-matchessuppliedexpandinguniversefigure}
  \lean{PhyXMiniProblems.ProblemPhyXMini0632.MatchesSuppliedExpandingUniverseFigure}
  \uses{def:physics:phyx-mini-0632:phyxminiproblems-problemphyxmini0632-society, def:physics:phyx-mini-0632:phyxminiproblems-problemphyxmini0632-figurequantitylabel, def:physics:phyx-mini-0632:phyxminiproblems-problemphyxmini0632-expandingsphericaluniversesetup}
  Direct qualitative evidence from the supplied raster image
\end{definition}
\begin{proof}
  This is the declared model definition of Matches Supplied Expanding Universe Figure.
\end{proof}
\begin{definition}[Has Physical Expanding Universe Parameters]
  \label{def:physics:phyx-mini-0632:phyxminiproblems-problemphyxmini0632-hasphysicalexpandinguniverseparameters}
  \lean{PhyXMiniProblems.ProblemPhyXMini0632.HasPhysicalExpandingUniverseParameters}
  \uses{def:physics:phyx-mini-0632:phyxminiproblems-problemphyxmini0632-timequantity, def:physics:phyx-mini-0632:phyxminiproblems-problemphyxmini0632-lengthinmeters, def:physics:phyx-mini-0632:phyxminiproblems-problemphyxmini0632-timeinseconds, def:physics:phyx-mini-0632:phyxminiproblems-problemphyxmini0632-speedinmeterspersecond, def:physics:phyx-mini-0632:phyxminiproblems-problemphyxmini0632-expandingsphericaluniversesetup}
  Positivity and minor-arc conditions selecting the physical branch shown
\end{definition}
\begin{proof}
  This is the declared model definition of Has Physical Expanding Universe Parameters.
\end{proof}
\begin{definition}[Satisfies Constant Expansion And Spherical Arc Laws]
  \label{def:physics:phyx-mini-0632:phyxminiproblems-problemphyxmini0632-satisfiesconstantexpansionandsphericalarclaws}
  \lean{PhyXMiniProblems.ProblemPhyXMini0632.SatisfiesConstantExpansionAndSphericalArcLaws}
  \uses{def:physics:phyx-mini-0632:phyxminiproblems-problemphyxmini0632-timequantity, def:physics:phyx-mini-0632:phyxminiproblems-problemphyxmini0632-lengthreadout, def:physics:phyx-mini-0632:phyxminiproblems-problemphyxmini0632-timereadout, def:physics:phyx-mini-0632:phyxminiproblems-problemphyxmini0632-speedreadout, def:physics:phyx-mini-0632:phyxminiproblems-problemphyxmini0632-expandingsphericaluniversesetup}
  The two governing laws are kept separate from the problem readouts: constant radial expansion gives R(t) = R₀ + v (t - t₀) in every compatible choice of length and time units; on a sphere, the minor surface arc between sites at fixed central angular separation has length D(t) = R(t) theta, with theta in radians. Neither law contains the requested numerical distance or an answer choice
\end{definition}
\begin{proof}
  This is the declared model definition of Satisfies Constant Expansion And Spherical Arc Laws.
\end{proof}
\begin{definition}[Answer Choice]
  \label{def:physics:phyx-mini-0632:phyxminiproblems-problemphyxmini0632-answerchoice}
  \lean{PhyXMiniProblems.ProblemPhyXMini0632.AnswerChoice}
  Labels of the four distance choices supplied with the problem
\end{definition}
\begin{proof}
  This is the declared model definition of Answer Choice.
\end{proof}
\begin{definition}[distance In Meters]
  \label{def:physics:phyx-mini-0632:phyxminiproblems-problemphyxmini0632-answerchoice-distanceinmeters}
  \lean{PhyXMiniProblems.ProblemPhyXMini0632.AnswerChoice.distanceInMeters}
  \uses{def:physics:phyx-mini-0632:phyxminiproblems-problemphyxmini0632-answerchoice}
  Numerical metre readout printed beside each answer label
\end{definition}
\begin{proof}
  This is the declared model definition of distance In Meters.
\end{proof}
\begin{definition}[recorded Dataset Answer]
  \label{def:physics:phyx-mini-0632:phyxminiproblems-problemphyxmini0632-recordeddatasetanswer}
  \lean{PhyXMiniProblems.ProblemPhyXMini0632.recordedDatasetAnswer}
  \uses{def:physics:phyx-mini-0632:phyxminiproblems-problemphyxmini0632-answerchoice}
  The answer label recorded in the source dataset
\end{definition}
\begin{proof}
  This is the declared model definition of recorded Dataset Answer.
\end{proof}
\begin{definition}[Rounds To Nearest Meter]
  \label{def:physics:phyx-mini-0632:phyxminiproblems-problemphyxmini0632-roundstonearestmeter}
  \lean{PhyXMiniProblems.ProblemPhyXMini0632.RoundsToNearestMeter}
  \uses{def:physics:phyx-mini-0632:phyxminiproblems-problemphyxmini0632-lengthquantity, def:physics:phyx-mini-0632:phyxminiproblems-problemphyxmini0632-lengthinmeters}
  A physical distance rounds to a displayed whole-metre value
\end{definition}
\begin{proof}
  This is the declared model definition of Rounds To Nearest Meter.
\end{proof}
\begin{definition}[Matches Answer Choice]
  \label{def:physics:phyx-mini-0632:phyxminiproblems-problemphyxmini0632-matchesanswerchoice}
  \lean{PhyXMiniProblems.ProblemPhyXMini0632.MatchesAnswerChoice}
  \uses{def:physics:phyx-mini-0632:phyxminiproblems-problemphyxmini0632-expandingsphericaluniversesetup, def:physics:phyx-mini-0632:phyxminiproblems-problemphyxmini0632-answerchoice, def:physics:phyx-mini-0632:phyxminiproblems-problemphyxmini0632-roundstonearestmeter}
  The observed surface distance rounds to the number printed for a choice
\end{definition}
\begin{proof}
  This is the declared model definition of Matches Answer Choice.
\end{proof}
\begin{definition}[Is Unique Matching Answer Choice]
  \label{def:physics:phyx-mini-0632:phyxminiproblems-problemphyxmini0632-isuniquematchinganswerchoice}
  \lean{PhyXMiniProblems.ProblemPhyXMini0632.IsUniqueMatchingAnswerChoice}
  \uses{def:physics:phyx-mini-0632:phyxminiproblems-problemphyxmini0632-expandingsphericaluniversesetup, def:physics:phyx-mini-0632:phyxminiproblems-problemphyxmini0632-answerchoice, def:physics:phyx-mini-0632:phyxminiproblems-problemphyxmini0632-matchesanswerchoice}
  A choice is the unique displayed value matching the rounded distance
\end{definition}
\begin{proof}
  This is the declared model definition of Is Unique Matching Answer Choice.
\end{proof}
\begin{theorem}[distance After Four Years exact readout]
  \label{thm:physics:phyx-mini-0632:phyxminiproblems-problemphyxmini0632-distanceafterfouryears-exact-readout}
  \lean{PhyXMiniProblems.ProblemPhyXMini0632.distanceAfterFourYears_exact_readout}
  \uses{def:physics:phyx-mini-0632:phyxminiproblems-problemphyxmini0632-lengthinmeters, def:physics:phyx-mini-0632:phyxminiproblems-problemphyxmini0632-fourproblemyearsinseconds, def:physics:phyx-mini-0632:phyxminiproblems-problemphyxmini0632-expandingsphericaluniversesetup, def:physics:phyx-mini-0632:phyxminiproblems-problemphyxmini0632-matchesexpandinguniverseproblemdata, def:physics:phyx-mini-0632:phyxminiproblems-problemphyxmini0632-hasphysicalexpandinguniverseparameters, def:physics:phyx-mini-0632:phyxminiproblems-problemphyxmini0632-satisfiesconstantexpansionandsphericalarclaws}
  Substitution of the calibrated data into the two governing laws gives the exact metre readout before rounding
\end{theorem}
\begin{proof}
  The conclusion follows from the governing relations and figure readouts named in the statement of distance After Four Years exact readout.
\end{proof}
% --- Archon declaration topology end ---
% --- Archon physics formalization source end ---
